\documentclass[a4paper,12pt]{article}
\usepackage[utf8]{inputenc}
\usepackage{geometry}
\usepackage{graphicx}
\usepackage{amsmath}
\usepackage{booktabs}
\usepackage{float}
\usepackage{hyperref}
\usepackage{listings}

% Page Margins
\geometry{left=2.54cm, right=2.54cm, top=2.54cm, bottom=2.54cm}

\begin{document}

\begin{titlepage}
    \centering
    \vspace*{\fill} % Pushes content down to the center vertically
    
    {\Huge \textbf{CSE 4512: Computer Networks Lab}\\ Lab 02 Report \par}
    
    \vspace{1.5cm} % Adds some space between title and author
    
    {\Large
        \textbf{Name: Rahinur Bin Naushad} \\
        Student ID: 220041118 \\
        Section: 1(B) \par
    }
    
    \vspace*{\fill} % Pushes the date to the bottom
    
    {\Large \today \par}
\end{titlepage}

\section{Lab Task}
The objective of this lab is to design a network topology for "Iutians Assemble," a company expanding its headquarters across five floors. The design requires Variable Length Subnet Masking (VLSM) to efficiently allocate IP addresses based on the specific device requirements of each department.

Given the Student ID \textbf{18}, the base network prefix is calculated as:
\[ 192.(18).(18+3 \pmod{14}) \rightarrow \textbf{192.18.7.0} \]

The device requirements are calculated based on the provided formulas:
\begin{itemize}
    \item \textbf{Engineering (Floor 1):} $E(18) = (7(18) + 3) \pmod{41} + 10 = 16$ devices.
    \item \textbf{Marketing (Floor 2):} $M(18) = (5(18) + 11) \pmod{26} + 5 = 28$ devices.
    \item \textbf{Sales (Floor 3):} $S(18) = (11(18) + 7) \pmod{36} + 5 = 30$ devices.
    \item \textbf{Administration (Floor 4):} $A(18) = (3(18) + 17) \pmod{18} + 3 = 20$ devices.
    \item \textbf{Management (Floor 5):} $G(18) = (13(18) + 1) \pmod{9} + 2 = 3$ devices.
\end{itemize}

\section{Network Topology}
Below is the screenshot of the final network topology designed in Cisco Packet Tracer.

\begin{figure}[H]
    \centering
    % REPLACE 'user' WITH YOUR ACTUAL USERNAME
    \includegraphics[width=0.9\textwidth]{/home/arcane/Pictures/Screenshots/topology.png}
    \caption{Final Network Topology for Iutians Assemble}
    \label{fig:topology}
\end{figure}

\section{Procedure}

\subsection{Step 1: VLSM Calculation}
We sorted the departments by the number of required host addresses (descending order) to ensure the most efficient use of the IP space. The base network is \texttt{192.18.7.0}.

\begin{enumerate}
    \item \textbf{Sales (30 hosts):} Requires 32 IPs (Block size). Mask: /27.
    \item \textbf{Marketing (28 hosts):} Requires 32 IPs. Mask: /27.
    \item \textbf{Administration (20 hosts):} Requires 32 IPs. Mask: /27.
    \item \textbf{Engineering (16 hosts):} Requires 32 IPs. Mask: /27.
    \item \textbf{Management (3 hosts):} Requires 8 IPs. Mask: /29.
\end{enumerate}

\subsection{Step 2: Subnetting Worksheet Configuration}
Using the calculations above, the subnet table was populated:

\begin{table}[H]
\centering
\begin{tabular}{|l|c|c|c|l|}
\hline
\textbf{Department} & \textbf{Dev} & \textbf{Network} & \textbf{Mask} & \textbf{Usable Host Range} \\ \hline
Sales & 30 & 192.18.7.0 & /27 & 192.18.7.1 -- 192.18.7.30 \\ \hline
Marketing & 28 & 192.18.7.32 & /27 & 192.18.7.33 -- 192.18.7.62\\ \hline
Administration & 20 & 192.18.7.64 & /27 & 192.18.7.65 -- 192.18.7.94 \\ \hline
Engineering & 16 & 192.18.7.96 & /27 & 192.18.7.97 -- 192.18.7.126 \\ \hline
Management & 3 & 192.18.7.128 & /29 & 192.18.7.129 -- 192.18.7.134 \\ \hline
\end{tabular}
\caption{Completed Subnetting Worksheet}
\end{table}

\subsection{Step 3: Topology Setup}
\begin{enumerate}
    \item Launched Cisco Packet Tracer.
    \item Placed 1 Router (Generic) to act as the central gateway.
    \item Placed 5 Switches, one for each department.
    \item Connected at least one PC to each switch to represent the end devices.
    \item Connected switches to the Router GigabitEthernet interfaces.
\end{enumerate}

\subsection{Step 4: Device Configuration}
\begin{enumerate}
    \item \textbf{Router Configuration:} Assigned the first usable IP of each subnet to the respective router interface (Default Gateway).
    \begin{itemize}
        \item Sales Gateway: 192.18.7.1
        \item Marketing Gateway: 192.18.7.33
        \item Admin Gateway: 192.18.7.65
        \item Engineering Gateway: 192.18.7.97
        \item Management Gateway: 192.18.7.129
    \end{itemize}
    \item \textbf{PC Configuration:} Assigned a usable IP to the PCs and set the Default Gateway to the router IP configured above. Subnet masks were set according to the  /27 or /29 CIDR values.
\end{enumerate}

\section{Observations \& Results}
To verify the network flow, ping tests were initiated from the \textbf{Sales Department PC} to the erst of the departments.

\begin{figure}[H]
    \centering
    \includegraphics[width=0.8\textwidth]{/home/arcane/Pictures/Screenshots/ping_result.png}
    \caption{Successful Ping from Sales (192.18.7.10) to Marketing (192.18.7.40)}
    \label{fig:ping}
\end{figure}

\begin{figure}[H]
    \centering
    \includegraphics[width=0.8\textwidth]{/home/arcane/Pictures/Screenshots/ping_result2.png}
    \caption{Successful Ping from Sales (192.18.7.10) to Administration (192.18.7.70)}
    \label{fig:ping}
\end{figure}

\begin{figure}[H]
    \centering
    \includegraphics[width=0.8\textwidth]{/home/arcane/Pictures/Screenshots/ping_result3.png}
    \caption{Successful Ping from Sales (192.18.7.10) to Engineering (192.18.7.100)}
    \label{fig:ping}
\end{figure}

\begin{figure}[H]
    \centering
    \includegraphics[width=0.8\textwidth]{/home/arcane/Pictures/Screenshots/ping_result4.png}
    \caption{Successful Ping from Sales (192.18.7.10) to Management (192.18.7.130)}
    \label{fig:ping}
\end{figure}


As seen in the screenshot, the packets were successfully transmitted and received, confirming that the routing table is correctly populated and the gateways are functioning.

\section{Challenges Faced}
\begin{itemize}
    \item \textbf{Router RIP Setup:} I set the router to router interface ip addresses in 192.168.0.0 network, which I needed to specify in the cli when using router rip command.
\end{itemize}

\end{document}
